The equivalent source method is widely used for processing and
interpolating magnetic data, particularly in airborne surveys. However,
implementations based on Cartesian coordinates present limitations at
regional and global scales, where Earth curvature introduces geometric
inconsistencies that affect data integration and modeling accuracy. To
address this problem, this study proposes an adaptation of the magnetic
equivalent source method to spherical coordinates, including revisions
to its mathematical formulation to account for spherical geometry. The
proposed framework enables consistent magnetic field modeling over large
geographic areas.
To improve the representation of magnetic sources, a dual-layer
configuration is adopted to separate long- and short-wavelength
components. Cross-validation is employed to determine optimal
hyperparameters for each layer, ensuring stable and balanced inversions.
To guarantee computational feasibility for large and high-resolution
datasets, a gradient-boosting strategy is incorporated into the inversion
process, significantly improving computational performance.
Synthetic experiments demonstrate that the method remains stable and accurate for large-scale datasets, 
with tests conducted on synthetic data containing up to 500,000 observations and enables
the reliable recovery of magnetic field components from total-field
anomaly data. The approach was further applied to more than 1.5 million
real observations, confirming its scalability and robustness. The
recovered field amplitude provides additional constraints for data
interpretation and enhances the geological analysis.
The final implementation is released as open-source software to support
reproducibility and broader adoption.